% s-349-p349-m002; original PDF 349, 350
\Aphorism{s-349-p349-m002}{400}
Causalität, Wirkung, Abhängigkeit: diese Begriffe werden immer noch zu menschlich, zu gegenständlich, nicht mathematisch genug verstanden. Auch der kälteste Verstand sieht etwas wie eine persönliche Beziehung zwischen Ursache und Wirkung, eine Art Gesetz und Kraftäusserung auf der einen, Gehorsam und Leiden auf der anderen Seite, während als Thatbestand nur das Zusammenauftreten Beider zu verzeichnen ist, ohne Activität und Passivität, ohne den Nebensinn, dass das Eine \Sperr{durch} das Andere dawäre. Dass es unlösbare Verknüpfungen giebt, ist das Wesentliche; welchen Richtungssinn man ihnen beilegt, ob den finalen oder effectiven, kommt nicht in Betracht. Eine Ursache, die zeitlich vorangeht, ist um nichts verständlicher als eine, die zeitlich folgt; weil es eben überhaupt weder Ursachen noch Wirkungen, Mittel noch Zwecke giebt, sondern nur Paare, inséparables von Dingen, deren eines nie ohne das andere auftritt. Der Missbrauch, den die Philosophie mit dem Begriffe Causalität treibt, pflanzt sich auf ein weiteres Gebiet, das der Naturwissenschaften, fort, als abergläubische Deutung des Begriffes „Kraft“. Man findet durch Erfahrung, dass die Bewegungen verschiedener materieller Punkte nicht unabhängig von einander sind; man erhält für diese Abhängigkeit eine besonders elegante Darstellung, wenn man neben den Variablen der Bewegung ihre ersten und zweiten Differentialquotienten nach der Zeit einführt, — eine Darstellung, die, als blosse analytische Zeichenschrift, die wirkliche Bewegung wiederum abzulesen gestattet. Aber was hat man für Mythologie mit diesem einfachen Sachverhalt getrieben! Die zweiten Differentialquotienten nennt man „Beschleunigungen“; sie würden verschwinden, wenn die ersten Differentialquotienten, die „Geschwindigkeiten“, constant wären. Änderungen der Geschwindigkeit sind aber Wirkungen einer „Kraft“; die materiellen Punkte, deren jeder eine vom Orte der anderen abhängige „Beschleunigung“ erleidet, üben also Kraftwirkung auf einander aus, sie „ziehen einander an“ oder „stossen einander ab“. Damit hat man den stupenden Begriff Gravitation gewonnen, der eine rein mathematische Beziehung zwischen verschiedenen Grössen in ein unbeholfen sinnlich gedachtes Schieben, Stossen, Drängen und Ziehen zwischen verschiedenen Erdenklümpchen, genannt Molecülen, vergröbert. Welcher Fetischismus! welches anthropomorphistische Unvermögen, über die ärmlichen Begriffe Wollen, Wirken, Handeln hinwegzusehen!
\par

